\section{预训练模型已经会推理吗：从 greedy decoding 到 CoT decoding}

如果中间轨迹有用，下一个问题不是立即微调，而是先检查预训练分布中是否已经存在这样的轨迹。课堂提出一个反直觉假设：模型并非“完全不会”，而可能只是 greedy decoding 过早选中了一个短、局部高概率但错误的分支。于是研究对象从参数能力转向解码策略。

\subsection{“不会”可能只是第一个候选不对}

讲义先展示当时的常见信念：预训练 LLM 必须经过提示工程或微调才能产生推理。下一页直接否定它，强调“all we need is decoding”。读这两页时要保留限定词：这不是说任何预训练模型都能解决任何任务，而是说\textbf{不能仅凭 top-1 贪心输出失败，就断言其分布中不存在正确轨迹}。

\lecturefigure{slide-06.jpg}{常见信念：不经 prompting 或 finetuning，预训练模型不能推理}{6}

\lecturefigure{slide-07.jpg}{课堂反命题：预训练分布可能已经包含推理轨迹}{7}

\teachervoice{00:06:20--00:08:30，老师把重点从“再训练一次模型”移到“先检查生成分布”。讲义提醒：latent capability 与 top-1 behavior 不是同一概念；失败输出可能来自搜索和排序，也可能来自分布中根本没有正确解，两者必须用候选分析区分。}

\subsection{苹果题：候选生成、错误排序与答案置信度}

下面三页把一次失败拆成“候选是否存在”和“排序是否正确”两个阶段。苹果题的 greedy 输出是 5，但候选集合中已经出现“父亲有 5 个，总数 $3+5=8$”的正确轨迹。页 8 先证明多看候选可以找到正确解；页 9 再否定一个诱人的捷径——按响应长度选答案，因为长文本也可能自信地犯错；页 10 最后提出按\textbf{最终答案部分的模型置信度}排序。读图时应追踪同一组候选在三页中如何被不同选择规则重新排序，而不是把三页当作重复截图。

\lecturefigure{slide-08.jpg}{Greedy 失败时，候选集合中仍可能存在正确推理}{8}

\lecturefigure{slide-09.jpg}{长度不是可靠的推理质量指标}{9}

\lecturefigure{slide-10.jpg}{按最终答案置信度选择候选，而非按轨迹长度}{10}

对候选 $y^{(k)}=(r^{(k)},a^{(k)})$，可定义答案区间的平均对数概率：
\[
s_k=\frac{1}{|a^{(k)}|}\sum_{j=1}^{|a^{(k)}|}
\log p_\theta\!\left(a^{(k)}_j\mid x,r^{(k)},a^{(k)}_{<j}\right),
\qquad k^*=\arg\max_k s_k.
\]
其中 $s_k$ 是第 $k$ 个候选答案的归一化置信度，$|a^{(k)}|$ 是答案 token 数，$k^*$ 是最终选择。只看最后一个 token 适合标签式答案；多 token 数字、程序或自由文本应使用明确的答案边界和长度归一化，否则短答案会得到不公平优势。

\begin{warningbox}{置信度不是正确性证明}
模型概率来自自身分布，可能系统性过度自信；答案抽取也可能把单位、符号或程序主体切错。置信度适合做候选排序信号，不应单独承担 hallucination 检测、上线门槛或安全证明。
\end{warningbox}

\subsection{Chain-of-Thought Decoding 的两步算法}

CoT decoding 的结构很简单：先超越 greedy，探索多个生成候选；再根据最终答案置信度选择。它与 CoT prompting 不同：前者修改\textbf{解码和排序}，后者修改\textbf{输入提示}。同一个模型可以不加示例地做 CoT decoding，也可以在 CoT prompt 下继续使用 greedy；两条轴必须分开。

\lecturefigure{slide-11.jpg}{Chain-of-Thought Decoding：扩展候选，再按答案置信度排序}{11}

\begin{lstlisting}[language=Python,caption={按答案置信度重排 CoT 候选的最小伪代码}]
def cot_decode(model, prompt, samples, answer_parser):
    candidates = model.sample(prompt, n=samples)
    ranked = []
    for response in candidates:
        answer_span = answer_parser(response.text)
        if answer_span is None:
            continue
        score = mean(response.logprobs[answer_span])
        ranked.append((score, response.text))
    if not ranked:
        raise ValueError("no parseable answer")
    return max(ranked)[1]
\end{lstlisting}

\begin{knowledgebox}{工程检查表：候选排序需要什么}
必须记录采样温度与 top-$p$、候选数、答案解析器版本、答案区间 log-prob、重复候选比例和无效解析率。否则“解码提升”可能只是多花 token、改变格式或意外泄漏答案边界，无法复现实验。
\end{knowledgebox}

\teachervoice{00:10:24--00:12:10，老师强调分数看 final answer 的条件概率，而不是把整条长轨迹的概率相乘后直接比较。课堂提示：轨迹越长，联合概率天然越小；不做答案边界和归一化会把长度偏差误当质量。}

\subsection{本章小结}

这一章把“预训练模型不会推理”拆成两个可检验命题：正确轨迹是否存在于候选分布，以及解码器是否能把它排到第一。CoT decoding 解决第二个问题，但它依赖答案解析和概率校准；若分布中根本缺少正确轨迹，就必须转向提示或训练。

